\section{The Method of Frobenius I}

\noindent
In this section we begin to study series solutions of a homogeneous linear second order differential
equation with a regular singular point at $x_{0}=0$, so it can be written as
\begin{equation}
  \label{eq:7.5.1}x^{2}A(x)y''+xB(x)y'+C(x)y=0,
\end{equation}
where $A$, $B$, $C$ are polynomials and $A(0)\ne0$.

We'll see that \eqref{eq:7.5.1} always has at least one solution of the form
\[
  y=x^{r}\sum_{n=0}^{\infty}a_{n}x^{n}
\]
where $a_{0}\ne0$ and $r$ is a suitably chosen number. The method we will use to find solutions of
this form and other forms that we'll encounter in the next two sections is called
\href{https://mathshistory.st-andrews.ac.uk/Biographies/Frobenius/}{\color{blue}
\it the method of Frobenius}, and we'll call them
\href{https://mathshistory.st-andrews.ac.uk/Biographies/Frobenius/}{\it Frobenius solutions\/}.

It can be shown that the power series $\sum_{n=0}^{\infty}a_{n}x^{n}$ in a Frobenius solution of
\eqref{eq:7.5.1} converges on some open interval $(-\rho,\rho)$, where $0<\rho\le\infty$. However,
since $x^{r}$ may be complex for negative $x$ or undefined if $x=0$, we'll consider solutions
defined for positive values of $x$. Easy modifications of our results yield solutions defined for
negative values of $x$. (Exercise~\ref{exer:7.5.54}).

We'll restrict our attention to the case where $A$, $B$, and $C$ are polynomials of degree not
greater than two, so \eqref{eq:7.5.1} becomes
\begin{equation}
  \label{eq:7.5.2}x^{2}(\alpha_{0}+\alpha_{1}x+\alpha_{2}x^{2})y''+x(\beta_{0}+\beta_{1}x+\beta_{2}x^{2})y' +(\gamma_{0}+\gamma_{1}x+\gamma_{2}x^{2})y=0,
\end{equation}
where $\alpha_{i}$, $\beta_{i}$, and $\gamma_{i}$ are real constants and $\alpha_{0}\ne0$. Most
equations that arise in applications can be written this way. Some examples are
\begin{eqnarray*}
  \alpha x^{2}y''+\beta xy'+\gamma y & = & 0\mbox{\quad(Euler's equation)},   \\
  x^{2}y''+xy'+(x^{2}-\nu^{2})y      & = & 0\mbox{\quad (Bessel's equation)}, \\
  \arraytext{and}                                                             \\
  xy''+(1-x)y'+\lambda y             & = & 0,\mbox{\quad\href{https://mathshistory.st-andrews.ac.uk/Biographies/Laguerre/}{(}{\color{blue}Laguerre's equation})},
\end{eqnarray*}
where we would multiply the last equation through by $x$ to put it in the form \eqref{eq:7.5.2}.
However, the method of Frobenius can be extended to the case where $A$, $B$, and $C$ are functions
that can be represented by power series in $x$ on some interval that contains zero, and
$A(0)\ne0$ (Exercises~\ref{exer:7.5.57} and \ref{exer:7.5.58}).

The next two theorems will enable us to develop systematic methods for finding Frobenius solutions
of \eqref{eq:7.5.2}.

\begin{theorem}
  \color{blue}
  \label{thmtype:7.5.1} Let
  \[
    Ly= x^{2}(\alpha_{0}+\alpha_{1}x+\alpha_{2}x^{2})y''+x(\beta_{0}+\beta_{1}x+\beta_{2}x^{2})y' +(\gamma_{0}+\gamma_{1}x+\gamma_{2}x^{2})y,
  \]
  and define
  \begin{eqnarray*}
    p_{0}(r)&=&\alpha_{0}r(r-1)+\beta_{0}r+\gamma_{0},\\[1\jot] p_{1}(r)&=&\alpha_{1}r(r-1)+\beta_{1}r+\gamma_{1},\\[1\jot] p_{2}(r)&=&\alpha_{2}r(r-1)+\beta_{2}r+\gamma_{2}.\\
  \end{eqnarray*}
  Suppose the series
  \begin{equation}
    \label{eq:7.5.3}y=\sum_{n=0}^{\infty}a_{n}x^{n+r}
  \end{equation}
  converges on $(0,\rho)$. Then
  \begin{equation}
    \label{eq:7.5.4}Ly=\sum_{n=0}^{\infty}b_{n}x^{n+r}
  \end{equation}
  on $(0,\rho)$, where
  \begin{equation}
    \label{eq:7.5.5}
    \begin{array}{ccl}
      b_{0} & = & p_{0}(r)a_{0},                                                       \\[1\jot]
      b_{1} & = & p_{0}(r+1)a_{1}+p_{1}(r)a_{0},                                       \\[1\jot]
      b_{n} & = & p_{0}(n+r)a_{n}+p_{1}(n+r-1)a_{n-1}+p_{2}(n+r-2)a_{n-2},\quad n\ge2.
    \end{array}
  \end{equation}
\end{theorem}

\proof We begin by showing that if $y$ is given by \eqref{eq:7.5.3} and $\alpha$, $\beta$, and
$\gamma$ are constants, then
\begin{equation}
  \label{eq:7.5.6}\alpha x^{2}y''+\beta xy'+\gamma y= \sum_{n=0}^{\infty}p(n+r)a_{n}x^{n+r},
\end{equation}
where
\[
  p(r)=\alpha r(r-1)+\beta r+\gamma.
\]
Differentiating \eqref{eq:7.5.3} twice yields
\begin{equation}
  \label{eq:7.5.7}y'=\sum_{n=0}^{\infty}(n+r)a_{n}x^{n+r-1}
\end{equation}
and
\begin{equation}
  \label{eq:7.5.8}y''=\sum_{n=0}^{\infty}(n+r)(n+r-1)a_{n}x^{n+r-2}.
\end{equation}
Multiplying \eqref{eq:7.5.7} by $x$ and \eqref{eq:7.5.8} by $x^{2}$ yields
\[
  xy'=\sum_{n=0}^{\infty}(n+r)a_{n}x^{n+r}
\]
and
\[
  x^{2}y''=\sum_{n=0}^{\infty}(n+r)(n+r-1)a_{n}x^{n+r}.
\]
Therefore
\begin{eqnarray*}
  \alpha x^{2}y''+\beta xy'+\gamma y & = & \sum_{n=0}^{\infty}\left[\alpha(n+r)(n+r-1)+\beta(n+r)+\gamma\right]a_{n}x^{n+r} \\[2\jot]
                                     & = & \sum_{n=0}^{\infty}p(n+r)a_{n}x^{n+r},
\end{eqnarray*}
which proves \eqref{eq:7.5.6}.

Multiplying \eqref{eq:7.5.6} by $x$ yields
\begin{equation}
  \label{eq:7.5.9}x(\alpha x^{2}y''+\beta xy'+\gamma y)=\sum_{n=0}^{\infty}p(n+r)a_{n}x^{n+r+1}= \sum_{n=1}^{\infty}p(n+r-1)a_{n-1}x^{n+r}.
\end{equation}
Multiplying \eqref{eq:7.5.6} by $x^{2}$ yields
\begin{equation}
  \label{eq:7.5.10}x^{2}(\alpha x^{2}y''+\beta xy'+\gamma y)=\sum_{n=0}^{\infty}p(n+r)a_{n}x^{n+r+2}= \sum_{n=2}^{\infty}p(n+r-2)a_{n-2}x^{n+r}.
\end{equation}

To use these results, we rewrite
\[
  Ly= x^{2}(\alpha_{0}+\alpha_{1}x+\alpha_{2}x^{2})y''+x(\beta_{0}+\beta_{1}x+\beta_{2}x^{2})y' +(\gamma_{0}+\gamma_{1}x+\gamma_{2}x^{2})y
\]
as
\begin{equation}
  \label{eq:7.5.11}
  \begin{array}{ccl}
    Ly & = & \left(\alpha_{0}x^{2}y''+\beta_{0}xy' +\gamma_{0}y\right)+ x\left(\alpha_{1}x^{2}y''+\beta_{1}xy'+\gamma_{1}y\right) \\[2\jot]
       &   & +\
x^{2}\left(\alpha_{2}x^{2}y''+\beta_{2}xy'+\gamma_{2}y\right).
  \end{array}
\end{equation}
From \eqref{eq:7.5.6} with $p=p_{0}$,
\[
  \alpha_{0}x^{2}y''+\beta_{0}xy'+\gamma_{0}y=\sum_{n=0}^{\infty}p_{0}(n+r)a_{n}x^{n+r}.
\]
From \eqref{eq:7.5.9} with $p=p_{1}$,
\[
  x\left(\alpha_{1}x^{2}y''+\beta_{1}xy'+\gamma_{1}y\right)=\sum_{n=1}^{\infty}p_{1}(n+r-1)a_{n-1}x^{n+r}.
\]
From \eqref{eq:7.5.10} with $p=p_{2}$,
\[
  x^{2}\left(\alpha_{2}x^{2}y''+\beta_{2}xy'+\gamma_{2}y\right)=\sum_{n=2}^{\infty}p_{2}(n+r-2)a_{n-2}x^{n+r}.
\]
Therefore we can rewrite \eqref{eq:7.5.11} as
\begin{eqnarray*}
  Ly=\sum_{n=0}^{\infty}p_{0}(n+r)a_{n}x^{n+r}+\sum_{n=1}^{\infty}p_{1}(n+r-1)a_{n-1}x^{n+r} \\[2\jot]
  +\sum_{n=2}^{\infty}p_{2}(n+r-2)a_{n-2}x^{n+r},
\end{eqnarray*}
or
\begin{eqnarray*}
  Ly & = & p_{0}(r)a_{0}x^{r}+\left[p_{0}(r+1)a_{1}+p_{1}(r)a_{0}\right]x^{r+1} \\[2\jot]
     &   & +\sum_{n=2}^{\infty}\left[p_{0}(n+r)a_{n}+p_{1}(n+r-1)a_{n-1}+p_{2}(n+r-2)a_{n-2}\right]x^{n+r},
\end{eqnarray*}
which implies \eqref{eq:7.5.4} with $\{b_{n}\}$ defined as in \eqref{eq:7.5.5}.

\begin{theorem}
  \color{blue}
  \label{thmtype:7.5.2} Let
  \[
    Ly= x^{2}(\alpha_{0}+\alpha_{1}x+\alpha_{2}x^{2})y''+x(\beta_{0}+\beta_{1}x+\beta_{2}x^{2})y' +(\gamma_{0}+\gamma_{1}x+\gamma_{2}x^{2})y,
  \]
  where $\alpha_{0}\ne0$, and define
  \begin{eqnarray*}
    p_{0}(r)&=&\alpha_{0}r(r-1)+\beta_{0}r+\gamma_{0},\\[1\jot] p_{1}(r)&=&\alpha_{1}r(r-1)+\beta_{1}r+\gamma_{1},\\[1\jot] p_{2}(r)&=&\alpha_{2}r(r-1)+\beta_{2}r+\gamma_{2}.\\
  \end{eqnarray*}
  Suppose $r$ is a real number such that $p_{0}(n+r)$ is nonzero for all positive integers $n$.
  Define
  \begin{equation}
    \label{eq:7.5.12}
    \begin{array}{ccl}
      a_{0}(r) & = & 1,                                                                                 \\
      a_{1}(r) & = & -\dst\frac{p_{1}(r)}{p_{0}(r+1)},                                                  \\[2\jot]
      a_{n}(r) & = & -\dst\frac{p_{1}(n+r-1)a_{n-1}(r)+p_{2}(n+r-2)a_{n-2}(r)}{p_{0}(n+r)},\quad n\ge2.
    \end{array}
  \end{equation}
  Then the Frobenius series
  \begin{equation}
    \label{eq:7.5.13}y(x,r)=x^{r}\sum_{n=0}^{\infty}a_{n}(r)x^{n}
  \end{equation}
  converges and satisfies
  \begin{equation}
    \label{eq:7.5.14}Ly(x,r)=p_{0}(r)x^{r}
  \end{equation}
  on the interval $(0,\rho)$, where $\rho$ is the distance from the origin to the nearest zero of
  $A(x)=\alpha_{0}+\alpha_{1}x+\alpha_{2}x^{2}$ in the complex plane. (If $A$ is constant, then
  $\rho=\infty$.)
\end{theorem}

If $\{a_{n}(r)\}$ is determined by the recurrence relation \eqref{eq:7.5.12} then substituting
$a_{n}=a_{n}(r)$ into \eqref{eq:7.5.5} yields $b_{0}=p_{0}(r)$ and $b_{n}=0$ for $n\ge1$, so
\eqref{eq:7.5.4} reduces to \eqref{eq:7.5.14}. We omit the proof that the series \eqref{eq:7.5.13}
converges on $(0,\rho)$.\bbox

If $\alpha_{i}=\beta_{i}=\gamma_{i}=0$ for $i=1$, $2$, then $Ly=0$ reduces to the Euler equation
\[
  \alpha_{0}x^{2}y''+\beta_{0}xy'+\gamma_{0}y=0.
\]
Theorem~\ref{thmtype:7.4.3} shows that the solutions of this equation are determined by the zeros of
the indicial polynomial
\[
  p_{0}(r)=\alpha_{0}r(r-1)+\beta_{0}r+\gamma_{0}.
\]
Since \eqref{eq:7.5.14} implies that this is also true for the solutions of $Ly=0$, we'll also say
that $p_{0}$ is the {\color{blue}\it indicial polynomial\/} of \eqref{eq:7.5.2}, and that
$p_{0}(r)=0$ is the {\color{blue}\it indicial equation\/} of $Ly=0$. We'll consider only cases where
the indicial equation has real roots $r_{1}$ and $r_{2}$, with $r_{1}\ge r_{2}$.

\begin{theorem}
  \color{blue}
  \label{thmtype:7.5.3} Let $L$ and $\{a_{n}(r)\}$ be as in Theorem~\ref{thmtype:7.5.2}, and
  suppose the indicial equation $p_{0}(r)=0$ of $Ly=0$ has real roots $r_{1}$ and $r_{2}$, where
  $r_{1}\ge r_{2}$. Then
  \[
    y_{1}(x)=y(x,r_{1})=x^{r_{1}}\sum_{n=0}^{\infty}a_{n}(r_{1})x^{n}
  \]
  is a Frobenius solution of $Ly=0$. Moreover, if $r_{1}-r_{2}$ isn't an integer then
  \[
    y_{2}(x)=y(x,r_{2})=x^{r_{2}}\sum_{n=0}^{\infty}a_{n}(r_{2})x^{n}
  \]
  is also a Frobenius solution of $Ly=0$, and $\{y_{1},y_{2}\}$ is a fundamental set of solutions.
\end{theorem}

\proof Since $r_{1}$ and $r_{2}$ are roots of $p_{0}(r)=0$, the indicial polynomial can be factored
as
\begin{equation}
  \label{eq:7.5.15}p_{0}(r)=\alpha_{0}(r-r_{1})(r-r_{2}).
\end{equation}
Therefore
\[
  p_{0}(n+r_{1})=n\alpha_{0}(n+r_{1}-r_{2}),
\]
which is nonzero if $n>0$, since $r_{1}-r_{2}\ge0$. Therefore the assumptions of
Theorem~\ref{thmtype:7.5.2} hold with $r=r_{1}$, and \eqref{eq:7.5.14} implies that
$Ly_{1}=p_{0}(r_{1})x^{r_{1}}=0$.

Now suppose $r_{1}-r_{2}$ isn't an integer. From \eqref{eq:7.5.15},
\[
  p_{0}(n+r_{2})=n\alpha_{0}(n-r_{1}+r_{2})\ne0\mbox{\quad if \quad}n=1,2,\cdots.
\]
Hence, the assumptions of Theorem~\ref{thmtype:7.5.2} hold with $r=r_{2}$, and \eqref{eq:7.5.14}
implies that $Ly_{2}=p_{0}(r_{2})x^{r_{2}}=0$. We leave the proof that $\{y_{1},y_{2}\}$ is a
fundamental set of solutions as an exercise (Exercise~\ref{exer:7.5.52}). \bbox

It isn't always possible to obtain explicit formulas for the coefficients in Frobenius solutions.
However, we can always set up the recurrence relations and use them to compute as many coefficients
as we want. The next example illustrates this.

\begin{example}
  \label{example:7.5.1}\rm Find a fundamental set of Frobenius solutions of
  \begin{equation}
    \label{eq:7.5.16}2x^{2}(1+x+x^{2})y''+x(9+11x+11x^{2})y'+(6+10x+7x^{2})y=0.
  \end{equation}
  Compute just the first six coefficients $a_{0}$,\dots, $a_{5}$ in each solution.
\end{example}

\solution For the given equation, the polynomials defined in Theorem~\ref{thmtype:7.5.2} are
\[
  \begin{array}{ccccc}
    p_{0}(r) & = & 2r(r-1)+9r+6   & = & (2r+3)(r+2), \\[1\jot]
    p_{1}(r) & = & 2r(r-1)+11r+10 & = & (2r+5)(r+2), \\[1\jot]
    p_{2}(r) & = & 2r(r-1)+11r+7  & = & (2r+7)(r+1).
  \end{array}
\]
The zeros of the indicial polynomial $p_{0}$ are $r_{1}=-3/2$ and $r_{2}=-2$, so $r_{1}-r_{2}=1/2$.
Therefore Theorem~\ref{thmtype:7.5.3} implies that
\begin{equation}
  \label{eq:7.5.17}y_{1}=x^{-3/2}\sum_{n=0}^{\infty}a_{n}(-3/2)x^{n}\quad\mbox{ and }\quad y_{2}=x^{-2}\sum_{n=0}^{\infty}a_{n}(-2)x^{n}
\end{equation}
form a fundamental set of Frobenius solutions of \eqref{eq:7.5.16}. To find the coefficients in
these series, we use the recurrence relation of Theorem~\ref{thmtype:7.5.2}; thus,
\begin{eqnarray*}
  a_{0}(r) & = & 1,                                                                                          \\
  a_{1}(r) & = & -\dst\frac{p_{1}(r)}{p_{0}(r+1)}=-\dst\frac{(2r+5)(r+2)}{(2r+5)(r+3)}=-\dst\frac{r+2}{r+3}, \\[2\jot]
  a_{n}(r) & = & -\dst\frac{p_{1}(n+r-1)a_{n-1}+p_{2}(n+r-2)a_{n-2}}{p_{0}(n+r)}                             \\[2\jot]
           & = & -\dst\frac{(n+r+1)(2n+2r+3)a_{n-1}(r)+(n+r-1)(2n+2r+3)a_{n-2}(r)}{(n+r+2)(2n+2r+3)}         \\[2\jot]
           & = & -\dst\frac{(n+r+1)a_{n-1}(r)+(n+r-1)a_{n-2}(r)}{n+r+2},\quad n\ge2.
\end{eqnarray*}

Setting $r=-3/2$ in these equations yields
\begin{equation}
  \label{eq:7.5.18}
  \begin{array}{ccl}
    a_{0}(-3/2) & = & 1,                                                                      \\
    a_{1}(-3/2) & = & -1/3,                                                                   \\[2\jot]
    a_{n}(-3/2) & = & -\dst\frac{(2n-1)a_{n-1}(-3/2)+ (2n-5)a_{n-2}(-3/2)}{2n+1},\quad n\ge2,
  \end{array}
\end{equation}
and setting $r=-2$ yields
\begin{equation}
  \label{eq:7.5.19}
  \begin{array}{ccl}
    a_{0}(-2) & = & 1,                                                            \\
    a_{1}(-2) & = & 0,                                                            \\[2\jot]
    a_{n}(-2) & = & -\dst\frac{(n-1)a_{n-1}(-2)+(n-3)a_{n-2}(-2)}{n},\quad n\ge2.
  \end{array}
\end{equation}
Calculating with \eqref{eq:7.5.18} and \eqref{eq:7.5.19} and substituting the results into
\eqref{eq:7.5.17} yields the fundamental set of Frobenius solutions
\begin{eqnarray*}
  y_{1} & = & x^{-3/2}\left(1-\frac{1}{3}x+\frac{2}{5}x^{2}-\frac{5}{21}x^{3}+\frac{7}{135}x^{4}+\frac{76}{1155}x^{5}+\cdots\right), \\[2\jot]
  y_{2} & = & x^{-2}\left(1+\frac{1}{2}x^{2}-\frac{1}{3}x^{3}+\frac{1}{8}x^{4}+\frac{1}{30}x^{5}+\cdots\right).
\end{eqnarray*}

\boxit{Special Cases With Two Term Recurrence Relations}

\noindent
For $n\ge2$, the recurrence relation \eqref{eq:7.5.12} of Theorem~\ref{thmtype:7.5.2} involves the
three coefficients $a_{n}(r)$, $a_{n-1}(r)$, and $a_{n-2}(r)$. We'll now consider some special cases
where \eqref{eq:7.5.12} reduces to a two term recurrence relation; that is, a relation involving
only $a_{n}(r)$ and $a_{n-1}(r)$ or only $a_{n}(r)$ and $a_{n-2}(r)$. This simplification often
makes it possible to obtain explicit formulas for the coefficients of Frobenius solutions.

We first consider equations of the form
\[
  x^{2}(\alpha_{0}+\alpha_{1}x)y''+x(\beta_{0}+\beta_{1}x)y'+(\gamma_{0}+\gamma_{1}x)y=0
\]
with $\alpha_{0}\ne0$. For this equation, $\alpha_{2}=\beta_{2}=\gamma_{2}=0$, so $p_{2}\equiv0$ and
the recurrence relations in Theorem~\ref{thmtype:7.5.2} simplify to
\begin{equation}
  \label{eq:7.5.20}
  \begin{array}{ccl}
    a_{0}(r) & = & 1,                                                          \\
    a_{n}(r) & = & -\dst\frac{p_{1}(n+r-1)}{p_{0}(n+r)}a_{n-1}(r),\quad n\ge1.
  \end{array}
\end{equation}

\begin{example}
  \label{example:7.5.2}\rm Find a fundamental set of Frobenius solutions of
  \begin{equation}
    \label{eq:7.5.21}x^{2}(3+x)y''+5x(1+x)y'-(1-4x)y=0.
  \end{equation}
  Give explicit formulas for the coefficients in the solutions.
\end{example}

\solution For this equation, the polynomials defined in Theorem~\ref{thmtype:7.5.2} are
\[
  \begin{array}{ccccc}
    p_{0}(r) & = & 3r(r-1)+5r-1 & = & (3r-1)(r+1), \\[1\jot]
    p_{1}(r) & = & r(r-1)+5r+4  & = & (r+2)^{2},   \\[1\jot]
    p_{2}(r) & = & 0.
  \end{array}
\]
The zeros of the indicial polynomial $p_{0}$ are $r_{1}=1/3$ and $r_{2}=-1$, so $r_{1}-r_{2}=4/3$.
Therefore Theorem~\ref{thmtype:7.5.3} implies that
\[
  y_{1}=x^{1/3}\sum_{n=0}^{\infty}a_{n}(1/3)x^{n}\quad\mbox{ and }\quad y_{2}=x^{-1}\sum_{n=0}^{\infty}a_{n}(-1)x^{n}
\]
form a fundamental set of Frobenius solutions of \eqref{eq:7.5.21}. To find the coefficients in
these series, we use the recurrence relations \eqref{eq:7.5.20}; thus,
\begin{equation}
  \label{eq:7.5.22}
  \begin{array}{ccl}
    a_{0}(r) & = & 1,                                                  \\
    a_{n}(r) & = & -\dst\frac{p_{1}(n+r-1)}{p_{0}(n+r)}a_{n-1}(r)      \\[2\jot]
             & = & -\dst\frac{(n+r+1)^{2}}{(3n+3r-1)(n+r+1)}a_{n-1}(r) \\[2\jot]
             & = & -\dst\frac{n+r+1}{3n+3r-1}a_{n-1}(r),\quad n\ge1.
  \end{array}
\end{equation}

Setting $r=1/3$ in \eqref{eq:7.5.22} yields
\begin{eqnarray*}
  a_{0}(1/3) & = & 1, \\
  a_{n}(1/3) & = & -\frac{3n+4}{9n}a_{n-1}(1/3),\quad n\ge1.
\end{eqnarray*}
By using the product notation introduced in Section~7.2 and proceeding as we did in the examples in
that section yields
\[
  a_{n}(1/3)=\frac{(-1)^{n}\prod_{j=1}^{n}(3j+4)}{9^{n}n!},\quad n\ge0.
\]
Therefore
\[
  y_{1}=x^{1/3}\sum_{n=0}^{\infty}\frac{(-1)^{n}\prod_{j=1}^{n}(3j+4)}{9^{n}n!}x^{n}
\]
is a Frobenius solution of \eqref{eq:7.5.21}.

Setting $r=-1$ in \eqref{eq:7.5.22} yields
\begin{eqnarray*}
  a_{0}(-1) & = & 1, \\
  a_{n}(-1) & = & -\frac{n}{3n-4}a_{n-1}(-1),\quad n\ge1,
\end{eqnarray*}
so
\[
  a_{n}(-1)=\frac{(-1)^{n}n!}{\prod_{j=1}^{n}(3j-4)}.
\]
Therefore
\[
  y_{2}=x^{-1}\sum_{n=0}^{\infty}\frac{(-1)^{n}n!}{\prod_{j=1}^{n}(3j-4)}x^{n}
\]
is a Frobenius solution of \eqref{eq:7.5.21}, and $\{y_{1},y_{2}\}$ is a fundamental set of
solutions. \bbox

We now consider equations of the form
\begin{equation}
  \label{eq:7.5.23}x^{2}(\alpha_{0}+\alpha_{2}x^{2})y''+x(\beta_{0}+\beta_{2}x^{2})y'+ (\gamma_{0}+\gamma_{2}x^{2})y=0
\end{equation}
with $\alpha_{0}\ne0$. For this equation, $\alpha_{1}=\beta_{1}=\gamma_{1}=0$, so $p_{1}\equiv0$ and
the recurrence relations in Theorem~\ref{thmtype:7.5.2} simplify to
\begin{eqnarray*}
  a_{0}(r) & = & 1, \\
  a_{1}(r) & = & 0, \\[2\jot]
  a_{n}(r) & = & -\dst\frac{p_{2}(n+r-2)}{p_{0}(n+r)}a_{n-2}(r),\quad n\ge2.
\end{eqnarray*}
Since $a_{1}(r)=0$, the last equation implies that $a_{n}(r)=0$ if $n$ is odd, so the Frobenius
solutions are of the form
\[
  y(x,r)=x^{r}\sum_{m=0}^{\infty}a_{2m}(r)x^{2m},
\]
where
\begin{equation}
  \label{eq:7.5.24}
  \begin{array}{ccl}
    a_{0}(r)  & = & 1,                                                             \\
    a_{2m}(r) & = & -\dst\frac{p_{2}(2m+r-2)}{p_{0}(2m+r)}a_{2m-2}(r),\quad m\ge1.
  \end{array}
\end{equation}

\begin{example}
  \label{example:7.5.3}\rm Find a fundamental set of Frobenius solutions of
  \begin{equation}
    \label{eq:7.5.25}x^{2}(2-x^{2})y''-x(3+4x^{2})y'+(2-2x^{2})y=0.
  \end{equation}
  Give explicit formulas for the coefficients in the solutions.
\end{example}

\solution For this equation, the polynomials defined in Theorem~\ref{thmtype:7.5.2} are
\[
  \begin{array}{ccccc}
    p_{0}(r) & = & 2r(r-1)-3r+2              & = & (r-2)(2r-1), \\[1\jot]
    p_{1}(r) & = & 0                          \\[1\jot]
    p_{2}(r) & = & -\left[r(r-1)+4r+2\right] & = & -(r+1)(r+2).
  \end{array}
\]
The zeros of the indicial polynomial $p_{0}$ are $r_{1}=2$ and $r_{2}=1/2$, so $r_{1}-r_{2}=3/2$.
Therefore Theorem~\ref{thmtype:7.5.3} implies that
\[
  y_{1}=x^{2}\sum_{m=0}^{\infty}a_{2m}(2)x^{2m}\quad\mbox{ and }\quad y_{2}=x^{1/2}\sum_{m=0}^{\infty}a_{2m}(1/2)x^{2m}
\]
form a fundamental set of Frobenius solutions of \eqref{eq:7.5.25}. To find the coefficients in
these series, we use the recurrence relation \eqref{eq:7.5.24}; thus,
\begin{equation}
  \label{eq:7.5.26}
  \begin{array}{ccl}
    a_{0}(r)  & = & 1,                                                                   \\
    a_{2m}(r) & = & -\dst\frac{p_{2}(2m+r-2)}{p_{0}(2m+r)}a_{2m-2}(r)                    \\[2\jot]
              & = & \dst\frac{(2m+r)(2m+r-1)}{(2m+r-2)(4m+2r-1)}a_{2m-2}(r),\quad m\ge1.
  \end{array}
\end{equation}

Setting $r=2$ in \eqref{eq:7.5.26} yields
\begin{eqnarray*}
  a_{0}(2)  & = & 1, \\
  a_{2m}(2) & = & \dst\frac{(m+1)(2m+1)}{m(4m+3)}a_{2m-2}(2),\quad m\ge1,
\end{eqnarray*}
so
\[
  a_{2m}(2)=(m+1)\prod_{j=1}^{m}\frac{2j+1}{4j+3}.
\]
Therefore
\enlargethispage{1in}
\[
  y_{1}=x^{2}\sum_{m=0}^{\infty}(m+1)\left(\prod_{j=1}^{m}\frac{2j+1}{4j+3}\right)x^{2m}
\]
is a Frobenius solution of \eqref{eq:7.5.25}.
\newpage

Setting $r=1/2$ in \eqref{eq:7.5.26} yields
\begin{eqnarray*}
  a_{0}(1/2)  & = & 1, \\
  a_{2m}(1/2) & = & \dst\frac{(4m-1)(4m+1)}{8m(4m-3)}a_{2m-2}(1/2),\quad m\ge1,
\end{eqnarray*}
so
\[
  a_{2m}(1/2)=\frac{1}{8^{m}m!}\prod_{j=1}^{m}\frac{(4j-1)(4j+1)}{4j-3}.
\]
Therefore
\[
  y_{2}=x^{1/2}\sum_{m=0}^{\infty}\frac{1}{8^{m}m!}\left(\prod_{j=1}^{m}\frac{(4j-1)(4j+1)}{4j-3}\right)x^{2m}
\]
is a Frobenius solution of \eqref{eq:7.5.25} and $\{y_{1},y_{2}\}$ is a fundamental set of
solutions.

\color{blue}
\remark{ Thus far, we considered only the case where the indicial equation has real roots that don't
differ by an integer, which allows us to apply Theorem~\ref{thmtype:7.5.3}. However, for equations
of the form \eqref{eq:7.5.23}, the sequence $\{a_{2m}(r)\}$ in \eqref{eq:7.5.24} is defined for
$r=r_{2}$ if $r_{1}-r_{2}$ isn't an {\color{blue}\it even\/} integer. It can be shown
(Exercise~\ref{exer:7.5.56}) that in this
case \[y_{1}=x^{r_{1}}\sum_{m=0}^{\infty}a_{2m}(r_{1})x^{2m}\quad\mbox{ and }\quad
y_{2}=x^{r_{2}}\sum_{m=0}^{\infty}a_{2m}(r_{2})x^{2m}\] form a fundamental set Frobenius solutions
of \eqref{eq:7.5.23}}.
\color{black}

\technology
\medskip

As we said at the end of Section~7.2, if you're interested in actually using series to compute
numerical approximations to solutions of a differential equation, then whether or not there's a
simple closed form for the coefficients is essentially irrelevant; recursive computation is usually
more efficient. Since it's also laborious, we encourage you to write short programs to implement
recurrence relations on a calculator or computer, even in exercises where this is not specifically
required.

In practical use of the method of Frobenius when $x_{0}=0$ is a regular singular point, we're
interested in how well the functions
\[
  y_{N}(x,r_{i})=x^{r_{i}}\sum_{n=0}^{N}a_{n}(r_{i})x^{n},\quad i=1,2,
\]
approximate solutions to a given equation when $r_{i}$ is a zero of the indicial polynomial. In
dealing with the corresponding problem for the case where $x_{0}=0$ is an ordinary point, we used
numerical integration to solve the differential equation subject to initial conditions
$y(0)=a_{0},\quad y'(0)=a_{1}$, and compared the result with values of the Taylor polynomial
\[
  T_{N}(x)=\sum_{n=0}^{N}a_{n}x^{n}.
\]
We can't do that here, since in general we can't prescribe arbitrary initial values for solutions of
a differential equation at a singular point. Therefore, motivated by Theorem~\ref{thmtype:7.5.2}
(specifically, \eqref{eq:7.5.14}), we suggest the following procedure.

\newpage
{\color{blue} \centerline{\bf Verification Procedure}

\it Let $L$ and $Y_{n}(x;r_{i})$ be defined by \[Ly=
x^{2}(\alpha_{0}+\alpha_{1}x+\alpha_{2}x^{2})y''+x(\beta_{0}+\beta_{1}x+\beta_{2}x^{2})y'
+(\gamma_{0}+\gamma_{1}x+\gamma_{2}x^{2})y\]
and \[y_{N}(x;r_{i})=x^{r_{i}}\sum_{n=0}^{N}a_{n}(r_{i})x^{n},\] where the coefficients
$\{a_{n}(r_{i})\}_{n=0}^{N}$ are computed as in \eqref{eq:7.5.12}, Theorem~\ref{thmtype:7.5.2}.
Compute the
error \begin{equation}\label{eq:7.5.27}E_{N}(x;r_{i})=x^{-r_{i}}Ly_{N}(x;r_{i})/\alpha_{0}\end{equation}
for various values of $N$ and various values of $x$ in the interval $(0,\rho)$, with $\rho$ as
defined in Theorem~\ref{thmtype:7.5.2}.}

\medskip

The multiplier $x^{-r_{i}}/\alpha_{0}$ on the right of \eqref{eq:7.5.27} eliminates the effects of
small or large values of $x^{r_{i}}$ near $x=0$, and of multiplication by an arbitrary constant. In
some exercises you will be asked to estimate the maximum value of $E_{N}(x;r_{i})$ on an interval
$(0,\delta]$ by computing $E_{N}(x_{m};r_{i})$ at the $M$ points $x_{m}=m\delta/M,\;m=1$, $2$,
\dots, $M$, and finding the maximum of the absolute values:
\begin{equation}
  \label{eq:7.5.28}\sigma_{N}(\delta)=\max\{|E_{N}(x_{m};r_{i})|,\;m=1,2,\dots,M\}.
\end{equation}
(For simplicity, this notation ignores the dependence of the right side of the equation on $i$ and
$M$.)

To implement this procedure, you'll have to write a computer program to calculate $\{a_{n}(r_{i})\}$
from the applicable recurrence relation, and to evaluate $E_{N}(x;r_{i})$.

The next exercise set contains five exercises specifically identified by \Lex that ask you to
implement the verification procedure. These particular exercises were chosen arbitrarily you can
just as well formulate such laboratory problems for any of the equations in any of the
Exercises~\ref{exer:7.5.1}--\ref{exer:7.5.10}, \ref{exer:7.5.14}-\ref{exer:7.5.25}, and
\ref{exer:7.5.28}--\ref{exer:7.5.51}

\exercises This set contains exercises specifically identified by \Lex that ask you to implement the verification procedure. These particular exercises were chosen arbitrarily you can just as well formulate such laboratory problems for any of the equations in Exercises~\ref{exer:7.5.1}--\ref{exer:7.5.10}, \ref{exer:7.5.14}-\ref{exer:7.5.25}, and \ref{exer:7.5.28}--\ref{exer:7.5.51}.

\noindent
{\it In Exercises~\ref{exer:7.5.1}--\ref{exer:7.5.10} find a fundamental set of Frobenius solutions.
Compute $a_{0}$, $a_{1}$ \dots, $a_{N}$ for $N$ at least $7$ in each solution.}
\begin{exerciselist}
  \item\label{exer:7.5.1} \Cex $2x^{2}(1+x+x^{2})y''+x(3+3x+5x^{2})y'-y=0$

  \item\label{exer:7.5.2} \Cex $3x^{2}y''+2x(1+x-2x^{2})y'+(2x-8x^{2})y=0$

  \item\label{exer:7.5.3} \Cex $x^{2}(3+3x+x^{2})y''+x(5+8x+7x^{2})y'-(1-2x-9x^{2})y=0$

  \item\label{exer:7.5.4} \Cex $4x^{2}y''+x(7+2x+4x^{2})y'-(1-4x-7x^{2})y=0$

  \item\label{exer:7.5.5} \Cex $12x^{2}(1+x)y''+x(11+35x+3x^{2})y'-(1-10x-5x^{2})y=0$

  \item\label{exer:7.5.6} \Cex $x^{2}(5+x+10x^{2})y''+x(4+3x+48x^{2})y'+(x+36x^{2})y=0$

  \item\label{exer:7.5.7} \Cex $8x^{2}y''-2x(3-4x-x^{2})y'+(3+6x+x^{2})y=0$

  \item\label{exer:7.5.8} \Cex $18x^{2}(1+x)y''+3x(5+11x+x^{2})y'-(1-2x-5x^{2})y=0$

  \item\label{exer:7.5.9} \Cex $x(3+x+x^{2})y''+(4+x-x^{2})y'+xy=0$

  \item\label{exer:7.5.10} \Cex $10x^{2}(1+x+2x^{2})y''+x(13+13x+66x^{2})y'-(1+4x+10x^{2})y=0$

  \item\label{exer:7.5.11} \Lex The Frobenius solutions of
  \[
    2x^{2}(1+x+x^{2})y''+x(9+11x+11x^{2})y'+(6+10x+7x^{2})y=0
  \]
  obtained in Example~\ref{example:7.5.1} are defined on $(0,\rho)$, where $\rho$ is defined in
  Theorem~\ref{thmtype:7.5.2}. Find $\rho$. Then do the following experiments for each Frobenius
  solution, with $M=20$ and $\delta=.5\rho$, $.7\rho$, and $.9\rho$ in the verification procedure
  described at the end of this section.
  \begin{alist}
    \item % (a)
    Compute $\sigma_{N}(\delta)$ (see Eqn.~\eqref{eq:7.5.28}) for $N=5$, $10$, $15$,\dots, $50$.

    \item % (b)
    Find $N$ such that $\sigma_{N}(\delta)<10^{-5}$.

    \item % (b)
    Find $N$ such that $\sigma_{N}(\delta)<10^{-10}$.
  \end{alist}

  \item\label{exer:7.5.12} \Lex By Theorem~\ref{thmtype:7.5.2} the Frobenius solutions of the
        equation in Exercise~\ref{exer:7.5.4} are defined on $(0,\infty)$. Do experiments \part{a},
        \part{b}, and \part{c} of Exercise~\ref{exer:7.5.11} for each Frobenius solution, with
        $M=20$ and $\delta=1$, $2$, and $3$ in the verification procedure described at the end of
        this section.

  \item\label{exer:7.5.13} \Lex The Frobenius solutions of the equation in Exercise~\ref{exer:7.5.6}
        are defined on $(0,\rho)$, where $\rho$ is defined in Theorem~\ref{thmtype:7.5.2}. Find
        $\rho$ and do experiments \part{a}, \part{b}, and \part{c} of Exercise~\ref{exer:7.5.11} for
        each Frobenius solution, with $M=20$ and $\delta=.3\rho$, $.4\rho$, and $.5\rho$, in the
        verification procedure described at the end of this section.

  \exercisetext{In Exercises~\ref{exer:7.5.14}--\ref{exer:7.5.25} find a fundamental set of
  Frobenius solutions. Give explicit formulas for the coefficients in each solution.}

  \item\label{exer:7.5.14} $2x^{2}y''+x(3+2x)y'-(1-x)y=0$

  \item\label{exer:7.5.15} $x^{2}(3+x)y''+x(5+4x)y'-(1-2x)y=0$

  \item\label{exer:7.5.16} $2x^{2}y''+x(5+x)y'-(2-3x)y=0$

  \item\label{exer:7.5.17} $3x^{2}y''+x(1+x)y'-y=0$

  \item\label{exer:7.5.18} $2x^{2}y''-xy'+(1-2x)y=0$

  \item\label{exer:7.5.19} $9x^{2}y''+9xy'-(1+3x)y=0$

  \item\label{exer:7.5.20} $3x^{2}y''+x(1+x)y'-(1+3x)y=0$

  \item\label{exer:7.5.21} $2x^{2}(3+x)y''+x(1+5x)y'+(1+x)y=0$

  \item\label{exer:7.5.22} $x^{2}(4+x)y''-x(1-3x)y'+y=0$

  \item\label{exer:7.5.23} $2x^{2}y''+5xy'+(1+x)y=0$

  \item\label{exer:7.5.24} $x^{2}(3+4x)y''+x(5+18x)y'-(1-12x)y=0$

  \item\label{exer:7.5.25} $6x^{2}y''+x(10-x)y'-(2+x)y=0$

  \item\label{exer:7.5.26} \Lex By Theorem~\ref{thmtype:7.5.2} the Frobenius solutions of the
        equation in Exercise~\ref{exer:7.5.17} are defined on $(0,\infty)$. Do experiments \part{a},
        \part{b}, and \part{c} of Exercise~\ref{exer:7.5.11} for each Frobenius solution, with
        $M=20$ and $\delta=3$, $6$, $9$, and $12$ in the verification procedure described at the end
        of this section.

  \item\label{exer:7.5.27} \Lex The Frobenius solutions of the equation in
        Exercise~\ref{exer:7.5.22} are defined on $(0,\rho)$, where $\rho$ is defined in
        Theorem~\ref{thmtype:7.5.2}. Find $\rho$ and do experiments \part{a}, \part{b}, and \part{c}
        of Exercise~\ref{exer:7.5.11} for each Frobenius solution, with $M=20$ and $\delta=.25\rho$,
        $.5\rho$, and $.75\rho$ in the verification procedure described at the end of this section.

  \newpage
  \exercisetext{In Exercises~\ref{exer:7.5.28}--\ref{exer:7.5.32} find a fundamental set of
  Frobenius solutions. Compute coefficients $a_{0}$, \dots, $a_{N}$ for $N$ at least $7$ in each
  solution.}

  \item\label{exer:7.5.28} \Cex $x^{2}(8+x)y''+x(2+3x)y'+(1+x)y=0$

  \item\label{exer:7.5.29} \Cex $x^{2}(3+4x)y''+x(11+4x)y'-(3+4x)y=0$

  \item\label{exer:7.5.30} \Cex $2x^{2}(2+3x)y''+x(4+11x)y'-(1-x)y=0$

  \item\label{exer:7.5.31} \Cex $x^{2}(2+x)y''+5x(1-x)y'-(2-8x)y=0$

  \item\label{exer:7.5.32} \Cex $x^{2}(6+x)y''+x(11+4x)y'+(1+2x)y=0$

  \exercisetext{In Exercises~\ref{exer:7.5.33}--\ref{exer:7.5.46} find a fundamental set of
  Frobenius solutions. Give explicit formulas for the coefficients in each solution.}

  \item\label{exer:7.5.33} $8x^{2}y''+x(2+x^{2})y'+y=0$

  \item\label{exer:7.5.34} $8x^{2}(1-x^{2})y''+2x(1-13x^{2})y'+(1-9x^{2})y=0$

  \item\label{exer:7.5.35} $x^{2}(1+x^{2})y''-2x(2-x^{2})y'+4y=0$

  \item\label{exer:7.5.36} $x(3+x^{2})y''+(2-x^{2})y'-8xy=0$

  \item\label{exer:7.5.37} $4x^{2}(1-x^{2})y''+x(7-19x^{2})y'-(1+14x^{2})y=0$

  \item\label{exer:7.5.38} $3x^{2}(2-x^{2})y''+x(1-11x^{2})y'+(1-5x^{2})y=0$

  \item\label{exer:7.5.39} $2x^{2}(2+x^{2})y''-x(12-7x^{2})y'+(7+3x^{2})y=0$

  \item\label{exer:7.5.40} $2x^{2}(2+x^{2})y''+x(4+7x^{2})y'-(1-3x^{2})y=0$

  \item\label{exer:7.5.41} $2x^{2}(1+2x^{2})y''+5x(1+6x^{2})y'-(2-40x^{2})y=0$

  \item\label{exer:7.5.42} $3x^{2}(1+x^{2})y''+5x(1+x^{2})y'-(1+5x^{2})y=0$

  \item\label{exer:7.5.43} $x(1+x^{2})y''+(4+7x^{2})y'+8xy=0$

  \item\label{exer:7.5.44} $x^{2}(2+x^{2})y''+x(3+x^{2})y'-y=0$

  \item\label{exer:7.5.45} $2x^{2}(1+x^{2})y''+x(3+8x^{2})y'-(3-4x^{2})y=0$

  \item\label{exer:7.5.46} $9x^{2}y''+3x(3+x^{2})y'-(1-5x^{2})y=0$

  \exercisetext{In Exercises~\ref{exer:7.5.47}--\ref{exer:7.5.51} find a fundamental set of
  Frobenius solutions. Compute the coefficients $a_{0}$, \dots, $a_{2M}$ for $M$ at least $7$ in
  each solution.}

  \item\label{exer:7.5.47} \Cex $6x^{2}y''+x(1+6x^{2})y'+(1+9x^{2})y=0$

  \item\label{exer:7.5.48} \Cex $x^{2}(8+x^{2})y''+7x(2+x^{2})y'-(2-9x^{2})y=0$

  \item\label{exer:7.5.49} \Cex $9x^{2}(1+x^{2})y''+3x(3+13x^{2})y'-(1-25x^{2})y=0$

  \item\label{exer:7.5.50} \Cex $4x^{2}(1+x^{2})y''+4x(1+6x^{2})y'-(1-25x^{2})y=0$

  \item\label{exer:7.5.51} \Cex $8x^{2}(1+2x^{2})y''+2x(5+34x^{2})y'-(1-30x^{2})y=0$

  \item\label{exer:7.5.52} Suppose $r_{1}>r_{2}$, $a_{0}=b_{0}=1$, and the Frobenius series
  \[
    y_{1}=x^{r_{1}}\sum_{n=0}^{\infty}a_{n}x^{n}\quad\mbox{ and }\quad y_{2}=x^{r_{2}}\sum_{n=0}^{\infty}b_{n}x^{n}
  \]
  both converge on an interval $(0,\rho)$.

  \newpage
  \begin{alist}
    \item % (a)
    Show that $y_{1}$ and $y_{2}$ are linearly independent on $(0,\rho)$. \hint{Show that if $c_{1}$
    and $c_{2}$ are constants such that $c_{1}y_{1}+c_{2}y_{2}\equiv0$ on $(0,\rho)$,
    then \[c_{1}x^{r_{1}-r_{2}}\sum_{n=0}^{\infty}a_{n}x^{n}+
    c_{2}\sum_{n=0}^{\infty}b_{n}x^{n}=0,\quad0<x<\rho.\] Then let $x\to0+$ to conclude that
    $c_{2}=0$.}

    \item % (b)
    Use the result of \part{a} to complete the proof of Theorem~\ref{thmtype:7.5.3}.
  \end{alist}

  \item\label{exer:7.5.53} The equation
  \setcounter{equation}{0}
  \begin{equation}
    \label{eq:7.5.{exer:7.5.53}A}x^{2}y''+xy'+(x^{2}-\nu^{2})y=0
  \end{equation}
  is \href{https://mathshistory.st-andrews.ac.uk/Biographies/Bessel/}{\color{blue}
  \it Bessel's equation of order~$\nu$\/}. (Here $\nu$ is a parameter, and this use of ``order''
  should not be confused with its usual use as in ``the order of the equation.'') The solutions of
  \eqref{eq:7.5.{exer:7.5.53}A} are
  \href{https://mathshistory.st-andrews.ac.uk/Biographies/Bessel/}{\color{blue}
  \it Bessel functions of order\/} $\nu$.
  \begin{alist}
    \item % (a)
    Assuming that $\nu$ isn't an integer, find a fundamental set of Frobenius solutions of
    \eqref{eq:7.5.{exer:7.5.53}A}.

    \item % (b)
    If $\nu=1/2$, the solutions of \eqref{eq:7.5.{exer:7.5.53}A} reduce to familiar elementary
    functions. Identify these functions.
  \end{alist}

  \item\label{exer:7.5.54}
  \begin{alist}
    \item % (a)
    Verify that
    \[
      \frac{d}{dx}\left(|x|^{r}x^{n}\right)=(n+r)|x|^{r}x^{n-1}\mbox{\quad and \quad}{\frac{d^{2}}{dx^{2}}}\left(|x|^{r}x^{n}\right)=(n+r)(n+r-1)|x|^{r}x^{n-2}
    \]
    if $x\ne0$.

    \item % (b)
    Let
    \[
      Ly= x^{2}(\alpha_{0}+\alpha_{1}x+\alpha_{2}x^{2})y''+x(\beta_{0}+\beta_{1}x+\beta_{2}x^{2})y' +(\gamma_{0}+\gamma_{1}x+\gamma_{2}x^{2})y=0.
    \]
    Show that if $x^{r}\sum_{n=0}^{\infty}a_{n}x^{n}$ is a solution of $Ly=0$ on $(0,\rho)$ then
    $|x|^{r}\sum_{n=0}^{\infty}a_{n}x^{n}$ is a solution on $(-\rho,0)$ and $(0,\rho)$.
  \end{alist}

  \item\label{exer:7.5.55}
  \begin{alist}
    \item % (a)
    Deduce from Eqn.~\eqref{eq:7.5.20} that
    \[
      a_{n}(r)=(-1)^{n}\prod_{j=1}^{n}\frac{p_{1}(j+r-1)}{p_{0}(j+r)}.
    \]

    \item % (b)
    Conclude that if $p_{0}(r)=\alpha_{0}(r-r_{1})(r-r_{2})$ where $r_{1}-r_{2}$ is not an integer,
    then
    \[
      y_{1}=x^{r_{1}}\sum_{n=0}^{\infty}a_{n}(r_{1})x^{n}\quad\mbox{ and }\quad y_{2}=x^{r_{2}}\sum_{n=0}^{\infty}a_{n}(r_{2})x^{n}
    \]
    form a fundamental set of Frobenius solutions of
    \[
      x^{2}(\alpha_{0}+\alpha_{1}x)y''+x(\beta_{0}+\beta_{1}x)y'+(\gamma_{0}+\gamma_{1}x)y=0.
    \]

    \item % (c)
    Show that if $p_{0}$ satisfies the hypotheses of \part{b} then
    \[
      y_{1}=x^{r_{1}}\sum_{n=0}^{\infty}\frac{(-1)^{n}}{n!\prod_{j=1}^{n}(j+r_{1}-r_{2})}\left(\frac{\gamma_{1}}{\alpha_{0}}\right)^{n}x^{n}
    \]
    and
    \[
      y_{2}=x^{r_{2}}\sum_{n=0}^{\infty}\frac{(-1)^{n}}{n!\prod_{j=1}^{n}(j+r_{2}-r_{1})}\left(\frac{\gamma_{1}}{\alpha_{0}}\right)^{n}x^{n}
    \]
    form a fundamental set of Frobenius solutions of
    \[
      \alpha_{0}x^{2}y''+\beta_{0}xy'+(\gamma_{0}+\gamma_{1}x)y=0.
    \]
  \end{alist}

  \item\label{exer:7.5.56} Let
  \[
    Ly=x^{2}(\alpha_{0}+\alpha_{2}x^{2})y''+x(\beta_{0}+\beta_{2}x^{2})y'+ (\gamma_{0}+\gamma_{2}x^{2})y=0
  \]
  and define
  \[
    p_{0}(r)=\alpha_{0}r(r-1)+\beta_{0}r+\gamma_{0}\quad\mbox{ and }\quad p_{2}(r)=\alpha_{2}r(r-1)+\beta_{2}r+\gamma_{2}.
  \]
  \begin{alist}
    \item % (a)
    Use Theorem~\ref{thmtype:7.5.2} to show that if
    \setcounter{equation}{0}
    \begin{equation}
      \label{eq:7.5.{exer:7.5.56}A}
      \begin{array}{rcl}
        a_{0}(r)                                      & = & 1,             \\
        p_{0}(2m+r)a_{2m}(r)+p_{2}(2m+r-2)a_{2m-2}(r) & = & 0,\quad m\ge1,
      \end{array}
    \end{equation}
    then the Frobenius series $y(x,r)=x^{r}\sum_{m=0}^{\infty}a_{2m}x^{2m}$ satisfies
    $Ly(x,r)=p_{0}(r)x^{r}$.

    \item % (b)
    Deduce from \eqref{eq:7.5.{exer:7.5.56}A} that if $p_{0}(2m+r)$ is nonzero for every positive
    integer $m$ then
    \[
      a_{2m}(r)=(-1)^{m}\prod_{j=1}^{m}\frac{p_{2}(2j+r-2)}{p_{0}(2j+r)}.
    \]

    \item % (c)
    Conclude that if $p_{0}(r)=\alpha_{0}(r-r_{1})(r-r_{2})$ where $r_{1}-r_{2}$ is not an even
    integer, then
    \[
      y_{1}=x^{r_{1}}\sum_{m=0}^{\infty}a_{2m}(r_{1})x^{2m}\quad\mbox{ and }\quad y_{2}=x^{r_{2}}\sum_{m=0}^{\infty}a_{2m}(r_{2})x^{2m}
    \]
    form a fundamental set of Frobenius solutions of $Ly=0$.

    \item % (d)
    Show that if $p_{0}$ satisfies the hypotheses of \part{c} then
    \[
      y_{1}=x^{r_{1}}\sum_{m=0}^{\infty}\frac{(-1)^{m}}{2^{m}m!\prod_{j=1}^{m}(2j+r_{1}-r_{2})}\left(\frac{\gamma_{2}}{\alpha_{0}}\right)^{m}x^{2m}
    \]
    and
    \[
      y_{2}=x^{r_{2}}\sum_{m=0}^{\infty}\frac{(-1)^{m}}{2^{m}m!\prod_{j=1}^{m}(2j+r_{2}-r_{1})}\left(\frac{\gamma_{2}}{\alpha_{0}}\right)^{m}x^{2m}
    \]
    form a fundamental set of Frobenius solutions of
    \[
      \alpha_{0}x^{2}y''+\beta_{0}xy'+(\gamma_{0}+\gamma_{2}x^{2})y=0.
    \]
  \end{alist}

  \item\label{exer:7.5.57} Let
  \[
    Ly=x^{2}q_{0}(x)y''+xq_{1}(x)y'+q_{2}(x)y,
  \]
  where
  \[
    q_{0}(x)=\sum_{j=0}^{\infty}\alpha_{j}x^{j},\quad q_{1}(x)=\sum_{j=0}^{\infty}\beta_{j}x^{j},\quad q_{2}(x)=\sum_{j=0}^{\infty}\gamma_{j}x^{j},
  \]
  and define
  \[
    p_{j}(r)=\alpha_{j}r(r-1)+\beta_{j}r+\gamma_{j},\quad j=0,1,\dots.
  \]
  Let $y=x^{r}\sum_{n=0}^{\infty}a_{n}x^{n}$. Show that
  \[
    Ly=x^{r}\sum_{n=0}^{\infty}b_{n}x^{n},
  \]
  where
  \[
    b_{n}=\sum_{j=0}^{n}p_{j}(n+r-j)a_{n-j}.
  \]

  \item\label{exer:7.5.58}
  \begin{alist}
    \item % (a)
    Let $L$ be as in Exercise~\ref{exer:7.5.57}, and suppose $p_{0}(n+r)\ne0$ for every positive
    integer $n$. Show that if
    \[
      y(x,r)=x^{r}\sum_{n=0}^{\infty}a_{n}(r)x^{n}
    \]
    where
    \begin{eqnarray*}
      a_{0}(r) & = & 1, \\
      a_{n}(r) & = & -\dst\frac{1}{p_{0}(n+r)}\sum_{j=1}^{n}p_{j}(n+r-j)a_{n-j}(r),\quad n\ge1,
    \end{eqnarray*}
    then
    \[
      Ly(x,r)=p_{0}(r)x^{r}.
    \]

    \item % (b)
    Conclude that if
    \[
      p_{0}(r)=\alpha_{0}(r-r_{1})(r-r_{2})
    \]
    where $r_{1}-r_{2}$ isn't an integer then $y_{1}=y(x,r_{1})$ and $y_{2}=y(x,r_{2})$ are
    solutions of $Ly=0$.
  \end{alist}

  \item\label{exer:7.5.59} Let
  \[
    Ly=x^{2}(\alpha_{0}+\alpha_{q}x^{q})y''+x(\beta_{0}+\beta_{q}x^{q})y'+ (\gamma_{0}+\gamma_{q}x^{q})y
  \]
  where $q$ is a positive integer, and define
  \[
    p_{0}(r)=\alpha_{0}r(r-1)+\beta_{0}r+\gamma_{0}\quad\mbox{ and }\quad p_{q}(r)=\alpha_{q}r(r-1)+\beta_{q}r+\gamma_{q}.
  \]
  \begin{alist}
    \item % (a)
    Show that if
    \[
      y(x,r)=x^{r}\sum_{m=0}^{\infty}a_{qm}(r)x^{qm}
    \]
    where
    \setcounter{equation}{0}
    \begin{equation}
      \label{eq:7.5.{exer:7.5.59}A}
      \begin{array}{rcl}
        a_{0}(r)  & = & 1,                                                                            \\
        a_{qm}(r) & = & -\dst\frac{p_{q}\left(q(m-1)+r\right)}{p_{0}(qm+r)}a_{q(m-1)}(r),\quad m\ge1,
      \end{array}
    \end{equation}
    then
    \[
      Ly(x,r)=p_{0}(r)x^{r}.
    \]

    \item % (b)
    Deduce from \eqref{eq:7.5.{exer:7.5.59}A} that
    \[
      a_{qm}(r)=(-1)^{m}\prod_{j=1}^{m}\frac{p_{q}\left(q(j-1)+r\right)}{p_{0}(qj+r)}.
    \]

    \item % (c)
    Conclude that if $p_{0}(r)=\alpha_{0}(r-r_{1})(r-r_{2})$ where $r_{1}-r_{2}$ is not an integer
    multiple of $q$, then
    \[
      y_{1}=x^{r_{1}}\sum_{m=0}^{\infty}a_{qm}(r_{1})x^{qm}\quad\mbox{ and }\quad y_{2}=x^{r_{2}}\sum_{m=0}^{\infty}a_{qm}(r_{2})x^{qm}
    \]
    form a fundamental set of Frobenius solutions of $Ly=0$.

    \item % (d)
    Show that if $p_{0}$ satisfies the hypotheses of \part{c} then
    \[
      y_{1}=x^{r_{1}}\sum_{m=0}^{\infty}\frac{(-1)^{m}}{q^{m}m!\prod_{j=1}^{m}(qj+r_{1}-r_{2})}\left(\frac{\gamma_{q}}{\alpha_{0}}\right)^{m}x^{qm}
    \]
    and
    \[
      y_{2}=x^{r_{2}}\sum_{m=0}^{\infty}\frac{(-1)^{m}}{q^{m}m!\prod_{j=1}^{m}(qj+r_{2}-r_{1})}\left(\frac{\gamma_{q}}{\alpha_{0}}\right)^{m}x^{qm}
    \]
    form a fundamental set of Frobenius solutions of
    \[
      \alpha_{0}x^{2}y''+\beta_{0}xy'+(\gamma_{0}+\gamma_{q}x^{q})y=0.
    \]
  \end{alist}

  \item\label{exer:7.5.60}
  \begin{alist}
    \item % (a)
    Suppose $\alpha_{0}$, $\alpha_{1}$, and $\alpha_{2}$ are real numbers with $\alpha_{0}\ne0$, and
    $\{a_{n}\}_{n=0}^{\infty}$ is defined by
    \[
      \alpha_{0}a_{1}+\alpha_{1}a_{0}=0
    \]
    and
    \[
      \alpha_{0}a_{n}+\alpha_{1}a_{n-1}+\alpha_{2}a_{n-2}=0,\quad n\ge2.
    \]
    Show that
    \[
      (\alpha_{0}+\alpha_{1}x+\alpha_{2}x^{2})\sum_{n=0}^{\infty}a_{n}x^{n}=\alpha_{0}a_{0},
    \]
    and infer that
    \[
      \sum_{n=0}^{\infty}a_{n}x^{n}=\frac{\alpha_{0}a_{0}}{\alpha_{0}+\alpha_{1}x+\alpha_{2}x^{2}}.
    \]

    \item % (b)
    With $\alpha_{0}$, $\alpha_{1}$, and $\alpha_{2}$ as in \part{a}, consider the equation
    \setcounter{equation}{0}
    \begin{equation}
      \label{eq:7.5.{exer:7.5.60}A}x^{2}(\alpha_{0}+\alpha_{1}x+\alpha_{2}x^{2})y''+x(\beta_{0}+\beta_{1}x+\beta_{2}x^{2})y'+ (\gamma_{0}+\gamma_{1}x+\gamma_{2}x^{2})y=0,
    \end{equation}
    and define
    \[
      p_{j}(r)=\alpha_{j}r(r-1)+\beta_{j}r+\gamma_{j},\quad j=0,1,2.
    \]
    Suppose
    \[
      \frac{p_{1}(r-1)}{p_{0}(r)}= \frac{\alpha_{1}}{\alpha_{0}},\qquad \frac{p_{2}(r-2)}{p_{0}(r)}= \frac{\alpha_{2}}{\alpha_{0}},
    \]
    and
    \[
      p_{0}(r)=\alpha_{0}(r-r_{1})(r-r_{2}),
    \]
    where $r_{1}>r_{2}$. Show that
    \[
      y_{1}=\frac{x^{r_{1}}}{\alpha_{0}+\alpha_{1}x+\alpha_{2}x^{2}}\quad\mbox{ and }\quad y_{2}=\frac{x^{r_{2}}}{\alpha_{0}+\alpha_{1}x+\alpha_{2}x^{2}}
    \]
    form a fundamental set of Frobenius solutions of \eqref{eq:7.5.{exer:7.5.60}A} on any interval
    $(0,\rho)$ on which $\alpha_{0}+\alpha_{1}x+\alpha_{2}x^{2}$ has no zeros.
  \end{alist}

  \exercisetext{In Exercises~\ref{exer:7.5.61}--\ref{exer:7.5.68} use the method suggested by
  Exercise~\ref{exer:7.5.60} to find the general solution on some interval $(0,\rho)$.}

  \item\label{exer:7.5.61} $2x^{2}(1+x)y''-x(1-3x)y'+y=0$

  \item\label{exer:7.5.62} $6x^{2}(1+2x^{2})y''+x(1+50x^{2})y'+(1+30x^{2})y=0$

  \item\label{exer:7.5.63} $28x^{2}(1-3x)y''-7x(5+9x)y'+7(2+9x)y=0$

  \item\label{exer:7.5.64} $9x^{2}(5+x)y''+9x(5+3x)y'-(5-8x)y=0$

  \item\label{exer:7.5.65} $8x^{2}(2-x^{2})y''+2x(10-21x^{2})y'-(2+35x^{2})y=0$

  \item\label{exer:7.5.66} $4x^{2}(1+3x+x^{2})y''-4x(1-3x-3x^{2})y'+3(1-x+x^{2})y=0$

  \item\label{exer:7.5.67} $3x^{2}(1+x)^{2}y''-x(1-10x-11x^{2})y'+(1+5x^{2})y=0$

  \item\label{exer:7.5.68} $4x^{2}(3+2x+x^{2})y''-x(3-14x-15x^{2})y'+(3+7x^{2})y=0$
\end{exerciselist}
